\section{The case \texorpdfstring{\(K^2=9\)}{K squared equals 9}}
\label{sec:rankone}

We now treat a minimal surface with $p_g=q=0$ and $K^2=9$, whose
intersection lattice has rank one. The argument will reduce its diagonal
coefficient to two numerical divisor types. These types occur with
nonnegative multiplicities on the given surface. On an auxiliary surface
with the same Chern numbers, the canonical-divisor type instead has a
negative virtual multiplicity. Comparing the two relations will leave
just one coefficient to evaluate, on a Hilbert scheme of length two.

Throughout this Section, $c_*$ means the coefficient chosen by the fixed
recursion of Proposition~\ref{prop:universal}. Its numerical selection
rule is the same on both surfaces; uniqueness of a diagonal decomposition
on either surface is not required.

\begin{proposition}\label{prop:rankone}
Let $X$ be a smooth minimal projective complex surface of general type
with $p_g(X)=q(X)=0$ and $K_X^2=9$.  Take
\begin{equation}\label{r1:data}
 L_1=L_2=0,\qquad D=3K_X,\qquad n=20,\qquad p=7.
\end{equation}
The coefficient $c$ in the Chow relation of
Theorem~\ref{thm:chow-relation} is strictly positive.
\end{proposition}

\subsection{Effective line classes and their multiplicities}

Write $K=K_X$.  The reductions in Section~\ref{sec:reductions} give
$b_2(X)=1$ and $\int_X c_2(T_X)=3$.  The exponential sequence, using
$H^1(X,\mathcal O_X)=H^2(X,\mathcal O_X)=0$, identifies
$\Pic(X)$ with $H^2(X,\mathbb Z)$.  Its quotient by torsion has a
positive unimodular intersection pairing of rank one.  Choose its
generator $G$ on the ray of an ample divisor.  Then
\[
 G^2=1,\qquad K\equiv 3G,
\]
where $\equiv$ denotes numerical equivalence.  Indeed $K^2=9$ gives the
coefficient $\pm3$, and the nef and big canonical class has positive
intersection with an ample divisor.  Every irreducible curve has a
positive coefficient on the same ray: its intersection with an ample
divisor is positive.  Consequently $K$ has positive intersection with
every curve, and $K^2>0$.  The Nakai--Moishezon criterion proves that
$K$ is ample.

For the data \eqref{r1:data}, Equation~\eqref{eq:vdim} gives
\[
 v=2\cdot20-K(3K)=13>9=p+2.
\]
Write $H\equiv aG$ and $T\equiv(9-a)G$.  Effectivity of both lines
implies $0\leq a\leq9$.  Their virtual divisor dimensions and total
point length, by Equation~\eqref{eq:length}, are
\begin{equation}\label{r1:numerical-classes}
 m_H=\frac{a(a-3)}2,\qquad
 m_T=\frac{(9-a)(6-a)}2,\qquad
 n_1+n_2=20-9a+a^2.
\end{equation}
Here and below an effective numerical class can represent several
distinct line bundles; their contributions must all be counted.
The complete numerical list is
\[
\begin{array}{c|rrrrrrrrrr}
 a&0&1&2&3&4&5&6&7&8&9\\ \hline
 m_H&0&-1&-1&0&2&5&9&14&20&27\\
 m_T&27&20&14&9&5&2&0&-1&-1&0\\
 n_1+n_2&20&12&6&2&0&0&2&6&12&20
\end{array}
\]
Proposition~\ref{prop:fixed-pg0} eliminates $a=1,2,7,8$ because one
of the two virtual divisor dimensions is negative.  At $a=4,5$ the
point lengths are both zero, so Proposition~\ref{prop:universal} gives
$c_*=0$.  Only the outer numerical type $(3K,0)$ and the middle
numerical type $(2K,K)$, with both orders included, can contribute to
$c$.

A numerically trivial effective line bundle is trivial: a nonzero
effective divisor has positive intersection with an ample divisor,
whereas a section with empty zero divisor trivializes its line bundle.
Hence the outer type consists of the single unordered pair $(3K,0)$.
The lines numerically equivalent to $K$ are precisely $K+\epsilon$,
where $\epsilon$ runs over $\Pic(X)_{\mathrm{tors}}$.  This torsion
group is finite, since $H^2(X,\mathbb Z)$ is finitely generated.  The
untwisted line $K$ has no section.  For $\epsilon\neq0$, the preceding
observation gives $H^0(X,-\epsilon)=0$, and therefore
\[
 h^2(X,K+\epsilon)=0,\qquad
 \chi(K+\epsilon)=1,\qquad
 h^0(X,K+\epsilon)=1+h^1(X,K+\epsilon)>0.
\]
Every such line is effective.  Its complement $2K-\epsilon$ is also
effective: $(2K-\epsilon)-K$ is ample, so Kodaira vanishing and
Riemann--Roch give $h^0(X,2K-\epsilon)=10$.  The same argument gives
$h^0(X,3K)=28$.  Thus all the required high lines occur.

Define
\begin{equation}\label{r1:torsion-count}
 \tau_X=\#\bigl(\Pic(X)_{\mathrm{tors}}\setminus\{0\}\bigr).
\end{equation}
The notation $\tau_X$ keeps this count distinct from the equivariant
parameter $t$.  Let $c_{\mathrm{out}}$ be the sum of the universal
coefficients for the two orders of the numerical pair $(3K,0)$,
including every distribution $n_1+n_2=20$.  Let
$c_{\mathrm{mid}}$ be the analogous sum for the two orders of
$(2K,K)$ and all distributions $n_1+n_2=2$.  These numbers are
defined by Equation~\eqref{rel:fiber-contribution} and
Proposition~\ref{prop:universal}, including when the indicated precise
line bundles are not effective on $X$.  Their numerical inputs are
$K^2=9$, $\int_Xc_2(T_X)=3$, $\chi(\mathcal O_X)=1$, and
\eqref{r1:data}.

For each nontrivial $\epsilon$, the two orders of
$(2K-\epsilon,K+\epsilon)$ have exactly these middle numerical inputs.
Each unordered pair is counted once by $\epsilon$, since its two
members have different numerical classes.  The linear systems need
not have the same actual dimension as $\epsilon$ varies; their
$H^1$ obstruction factors have already been incorporated in the
coefficient extraction of Proposition~\ref{prop:fixed-pg0}.
It follows that the total coefficient on $X$ is
\begin{equation}\label{r1:target-coefficient}
 c=c_{\mathrm{out}}+\tau_Xc_{\mathrm{mid}}.
\end{equation}

The outer sum involves Hilbert schemes of total length $20$, whereas the
middle sum has total length $2$. We will use the auxiliary relation to
express $c_{\mathrm{out}}$ as a positive multiple of $c_{\mathrm{mid}}$,
avoiding the length-$20$ calculation. Positivity will then follow from
the explicit length-two computation.

\subsection{Comparison on the Cartwright--Steger surface}

Let $Y$ be the Cartwright--Steger surface.  Its existence as a smooth
projective compact two-ball quotient follows from the torsion-free
lattice in \cite[Theorem~2]{CKY}.  Its invariants are
\begin{equation}\label{r1:auxiliary-invariants}
 K_Y^2=9,\qquad \int_Yc_2(T_Y)=3,\qquad
 \chi(\mathcal O_Y)=1,\qquad q(Y)=p_g(Y)=1;
\end{equation}
see \cite[Main Theorem~(a) and Lemma~3]{CKY}.
The Bergman metric on the ball descends to $Y$, and its negative
Ricci form gives a positive curvature form on $K_Y$.  Kodaira's
positivity criterion therefore makes $K_Y$ ample.  Thus $Y$ satisfies
all the hypotheses of the auxiliary construction of
Proposition~\ref{prop:fixed-auxiliary} with $D=3K_Y$.

Choose a nonzero $\omega\in H^0(Y,K_Y)$ and set
$\alpha=[\omega+\overline\omega]\in H^2(Y,\mathbb R)$.  Holomorphic
top forms are closed, and
\begin{equation}\label{r1:test-class}
 \alpha^2:=\int_Y\alpha\smile\alpha
   =2\int_Y\omega\wedge\overline\omega>0.
\end{equation}
The inequality follows in local holomorphic coordinates from positivity
of the associated real volume form.  Hodge type gives
$\alpha\cdot P=0$ for every divisor class $P$ on $Y$.

\begin{lemma}\label{r1:auxiliary-comparison}
For the common universal coefficients defined above, one has
\begin{equation}\label{r1:comparison}
 c_{\mathrm{out}}-s c_{\mathrm{mid}}=0,
 \qquad s>0,
\end{equation}
where $-s$ is the canonical-divisor coefficient in
Equation~\eqref{eq:canonical-cycle} for $Y$.
\end{lemma}

\begin{proof}
Apply the construction on $Y$ with the data \eqref{r1:data}.
Proposition~\ref{prop:fixed-auxiliary} first replaces the low divisor
virtual class by the empty-divisor point with coefficient $1$, or by
the canonical-divisor point with coefficient $-s$. The compatible virtual
pullback along the high morphism is its smooth pullback multiplied by
the high point-obstruction Euler class. Thus the surviving terms are computed on
\[
 |3K_Y|\times Y^{[n_1]}\times Y^{[n_2]}
 \quad\text{or}\quad
 |2K_Y|\times Y^{[n_1]}\times Y^{[n_2]},
\]
with total point lengths $20$ and $2$, respectively, and with both fixed
point-obstruction Euler factors retained. The high linear systems have
dimensions $27$ and $9$. The low divisor is fixed, so its hyperplane
variable is zero. Both orders occur, and the canonical multiplier is
$-s$ in each order.

The projection formula restricts all operational insertions to these
actual fibers. Derived base change likewise restricts the moving
complexes, giving exactly Equation~\eqref{eq:moving-ratio}.
Consequently the remaining formulas are those defining
$c_{\mathrm{out}}$ and $c_{\mathrm{mid}}$, with their indicated
multipliers. This restriction precedes the use of universality: no
integration over a varying Picard component is being compared by
numerical invariants.

After that restriction, the equality of the numerical data in
\eqref{r1:auxiliary-invariants} with those for $X$ permits
Proposition~\ref{prop:universal} to identify the chosen coefficients.
The resulting Chow relation on $Y^2$ has diagonal coefficient
$c_{\mathrm{out}}-s c_{\mathrm{mid}}$.  Pair it with
$\alpha\boxtimes\alpha$.  Every external term vanishes: bidegree
allows only a divisor on each factor to pair nontrivially, and each
such divisor is orthogonal to $\alpha$.  The diagonal pairs to
$\alpha^2$.  Equation~\eqref{r1:test-class} now gives
\eqref{r1:comparison}.
\end{proof}

\subsection{The middle coefficient}

\begin{lemma}\label{lem:middle-coefficient}
For the data \eqref{r1:data}, with numerical inputs
$K^2=9$ and $\int c_2=3$, the two-order middle coefficient is
\[
 c_{\mathrm{mid}}=12.
\]
\end{lemma}

\begin{proof}
Compute on $Y$ with the class $\alpha$ in \eqref{r1:test-class},
omitting the overall canonical multiplier $-s$. We compute the order
$T=2K_Y$, $H=K_Y$ first, and verify equality with the reverse order
at the end. For this order,
\[
 n_1+n_2=2,\qquad m_T=9,\qquad m_H=0.
\]
The three possible point-length distributions are $(2,0)$, $(1,1)$,
and $(0,2)$.  Taking the coefficient of $x_2^0$ means setting
$x_2=0$: every moving weight in
Equation~\eqref{eq:moving-ratio} has nonzero constant term $1$ or
$-1$ after $t=1/2$, so all its inverse Euler factors expand with
nonnegative powers of $x_1,x_2$.  This remains true for virtual
bundles, by expanding their Chern classes at the nonzero weight.
If $n_2>0$, the low fixed factor at $x_2=0$ is
\[
 e\bigl((\mathcal O_Y^{[n_2]})^\vee\bigr)=0.
\]
Indeed the constant function $1$ is a nowhere-zero section of the
rank-$n_2$ tautological bundle, so its top Chern class, and hence
that of its dual, vanishes.  Equivalently its Euler polynomial is
divisible by $x_2$; there is no negative power of $x_2$ in the other
factors to cancel it.  Therefore only $(n_1,n_2)=(2,0)$ contributes.

Put
\[
 U=(-K_Y)^{[2]},\qquad N=\mathcal O_Y^{[2]},\qquad x=x_1,
 \qquad
 \mu=\int_Y[Z_1]\alpha\in H^2(Y^{[2]},\mathbb R).
\]
All the other terms of Equation~\eqref{eq:insertion} integrate to
zero against $\alpha$, by its orthogonality to divisors and by
degree.  Consequently $\mu$ is the class used in the pairing
formula \eqref{rel:pairing}.

For this order, the line data in the moving ratio are
\[
 P_{12}=2K_Y,\qquad P_{21}=K_Y,\qquad
 B_{12}=K_Y,\qquad B_{21}=-K_Y.
\]
Riemann--Roch gives
$\chi(K_Y)=1$ and $\chi(2K_Y)=\chi(-K_Y)=10$.  The scalar part of
the moving ratio is therefore
\begin{equation}\label{r1:scalar-ratio}
 \frac{(x-1)^1}{(x-1)^{10}}
 \frac{(1-x)^{10}}{1^1}=-(1-x).
\end{equation}
Since the second ideal is $\mathcal O_Y$, its two point-complex
terms are
\begin{equation}\label{r1:point-complexes}
 E_{12}(K_Y)=N^\vee,\qquad E_{21}(-K_Y)=U
 \quad\text{in }K(Y^{[2]}).
\end{equation}
For the first identity, apply $\RHom(-,K_Y)$ to the point-ideal
sequence: the difference
$[\RGamma(K_Y)]-[\RHom(\mathcal I_1,K_Y)]$ is
$[\RHom(\mathcal O_{Z_1},K_Y)]=[N^\vee[-2]]=[N^\vee]$.
For the second, apply derived sections to that sequence tensored
by $-K_Y$.  The even shift in the first identity contributes no
minus sign in $K$-theory.

For any rank-two bundle $F$ one has
$e_{-u}(F^\vee)=e_u(F)$.  Hence the moving point numerator
$e_{x-1}(U^\vee)$ cancels the denominator $e_{1-x}(U)$, and the
remaining denominator $e_{x-1}(N^\vee)$ equals $e_{1-x}(N)$.
Including the high fixed Euler factor $e_x(U^\vee)$ and the
insertion $x^7$, the pairing for this order is
\begin{equation}\label{r1:middle-integral}
 I=-[x^2]\int_{Y^{[2]}}
       \mu^2(1-x)\frac{e_x(U^\vee)}{e_{1-x}(N)}.
\end{equation}

Here is the required homogeneous expansion.  Set
$u_i=c_i(U)$ and $\nu_i=c_i(N)$.  Then
\[
 e_x(U^\vee)=x^2-u_1x+u_2,\qquad
 e_{1-x}(N)=(1-x)^2+\nu_1(1-x)+\nu_2,
\]
and, through degree two in the Chern classes on $Y^{[2]}$,
\[
 \frac{1-x}{e_{1-x}(N)}
 =\frac1{1-x}-\frac{\nu_1}{(1-x)^2}
       +\frac{\nu_1^2-\nu_2}{(1-x)^3}
       +\text{terms of Chern degree at least three}.
\]
The class $\mu^2$ has real cohomological degree four and $Y^{[2]}$
has complex dimension four.  Thus only Chern degree two in the
remaining factor can contribute to the integral.  Multiplication
by $x^2-u_1x+u_2$ and extraction of $x^2$ in that degree gives
$u_2+2u_1\nu_1+\nu_1^2-\nu_2$.  Equation~\eqref{r1:middle-integral}
therefore becomes
\begin{equation}\label{r1:chern-integral}
 I=-\int_{Y^{[2]}}\mu^2
   \bigl(c_2(U)+2c_1(U)c_1(N)+c_1(N)^2-c_2(N)\bigr).
\end{equation}

We have reduced the middle coefficient to four tautological Chern
terms in \eqref{r1:chern-integral}. To evaluate them, consider the
consecutive nested Hilbert scheme
\[
 \widetilde{Y^2}=Y^{[1,2]}=\operatorname{Bl}_{\Delta_Y}(Y^2),
 \qquad
 \pi:\widetilde{Y^2}\longrightarrow Y^2,
 \qquad
 \psi:\widetilde{Y^2}\longrightarrow Y^{[2]}.
\]
The first surface coordinate records the length-one subscheme and
the second the residual point.  The morphism $\psi$ is proper and
generically of degree two, since either point of a reduced length-two
subscheme can be selected as its length-one subscheme.  The
projectivization description and the tautological sequence are those
of \cite[\S1, equations~(3)--(6), and Lemma~2.1]{EGL}, specialized
to length one.  Write $E$ for the exceptional divisor and
$p_i=\operatorname{pr}_i\circ\pi$.  The tautological quotient line
of the projectivized diagonal ideal is $\mathcal O(-E)$.
Consequently, for any line bundle $P$ on $Y$, the tautological
sequence is
\begin{equation}\label{r1:tautological-sequence}
 0\longrightarrow p_2^*P\otimes\mathcal O(-E)
 \longrightarrow\psi^*(P^{[2]})
 \longrightarrow p_1^*P\longrightarrow0.
\end{equation}
It follows either from the universal point-ideal sequence, by
tensoring with $P$ and pushing along the finite universal subscheme,
or directly from the quotient recording the residual point.

Use subscripts for pullbacks by $p_i$, and set $P=-K_Y$.  The
length-two incidence cycle pulls back to the sum of the two point
graphs, so
\begin{equation}\label{r1:blowup-classes}
 \begin{gathered}
 \psi^*\mu=\alpha_1+\alpha_2,\qquad
 \psi^*c_1(N)=-E,\qquad \psi^*c_2(N)=0,\\
 \psi^*c_1(U)=P_1+P_2-E,\qquad
 \psi^*c_2(U)=P_1P_2-P_1E.
 \end{gathered}
\end{equation}
The incidence assertion also follows by taking the codimension-two
Chern character of the universal point-ideal sequence: its
$\mathcal O(-E)$ twist does not alter the graph's fundamental
codimension-two term.  The bundle identities in
\eqref{r1:blowup-classes} follow at once from
\eqref{r1:tautological-sequence}, including the sign of $E$.

For completeness, the exceptional-divisor pushforwards needed here are
\begin{equation}\label{r1:exceptional-push}
 \pi_*E=0,\qquad \pi_*(E^2)=-[\Delta_Y].
\end{equation}
The first vanishes because $E$ maps to a smaller-dimensional cycle.
For the second, $E$ is a $\mathbb P^1$-bundle over $\Delta_Y$ and
$\mathcal O(E)|_E=\mathcal O_E(-1)$, whose degree on each fiber is
$-1$.  This proves the sign directly.

The bracket in Equation~\eqref{r1:chern-integral} now pulls back to
\[
 P_1P_2-(3P_1+2P_2)E+3E^2.
\]
All terms linear in $E$ integrate to zero by
\eqref{r1:exceptional-push} and the projection formula.  The remaining
term involving $P$ satisfies
\[
 \int_{Y^2}(\alpha_1+\alpha_2)^2P_1P_2
       =2(\alpha\cdot P)^2=0;
\]
the two other terms in the square exceed the dimension of one
surface factor.  Since $\psi_*[\widetilde{Y^2}]=2[Y^{[2]}]$, the
factor of two from this cover must also be included.  Thus
\begin{align*}
 I
 &=-\frac12\int_{\widetilde{Y^2}}
                  (\alpha_1+\alpha_2)^2\,3E^2\\
 &=\frac32\int_{\Delta_Y}(\alpha_1+\alpha_2)^2
   =\frac32\int_Y(2\alpha)^2
   =6\alpha^2.
\end{align*}
The external terms pair to zero, so Equation~\eqref{rel:pairing}
and $\alpha^2\neq0$ show that this order has universal diagonal
coefficient $6$.

Finally exchange the two orders.  Since $L_1=L_2=0$, exchanging
the two summands is the same as replacing the equivariant parameter
$t$ by $-t$ before evaluation.  The ordinary codimension-two output
component has parameter power
\[
 p+2-v=7+2-13=-4.
\]
This exponent is even, so the exchange preserves its pairing at
$t=1/2$.  The second order therefore also has coefficient $6$.
Both orders and all three point-length distributions have now been
accounted for, and $c_{\mathrm{mid}}=6+6=12$.
\end{proof}

\begin{proof}[Proof of Proposition~\ref{prop:rankone}]
Combine Equations~\eqref{r1:target-coefficient} and
\eqref{r1:comparison} with Lemma~\ref{lem:middle-coefficient}.  On
the original surface $X$ they give
\[
 c=(s+\tau_X)c_{\mathrm{mid}}=12(s+\tau_X)>0,
\]
because $s>0$ and $\tau_X\geq0$.
\end{proof}
